\notepage{10}{Finding and normalizing eigenvectors}
For a known eigenvalue $\omega_k$, solve the homogeneous linear system
\[
 \begin{bmatrix}
 \Omega_{11}-\omega_k&\Omega_{12}&\cdots&\Omega_{1n}\\
 \Omega_{21}&\Omega_{22}-\omega_k&\cdots&\Omega_{2n}\\
 \vdots&\vdots&\ddots&\vdots\\
 \Omega_{n1}&\Omega_{n2}&\cdots&\Omega_{nn}-\omega_k
 \end{bmatrix}
 \begin{bmatrix}v_1\\v_2\\\vdots\\v_n\end{bmatrix}=0.
\]
For a one-dimensional eigenspace, this fixes the component ratios;
some components may be zero. The notes suggest checking the $n=2$
and $n=3$ cases directly.

\smallheading{Normalization convention}
The solutions $\{v_i\}$ and $\{\alpha v_i\}$ describe proportional
eigenvectors, $\ket{v}$ and $\alpha\ket{v}$. Their lengths can differ,
so we conventionally choose
\[
 \braket{v}{v}=1.
\]
Starting with any nonzero eigenvector, the normalized vector is
\[
 \boxed{\ket{v}_{\mathrm{norm}}
       =\frac{\ket{v}}{\sqrt{\braket{v}{v}}}.}
\]
Normalization leaves an overall phase undetermined.

\smallheading{Repeated eigenvalues}
If several eigenvalues coincide,
\[
 \omega_{k_1}=\omega_{k_2}=\cdots=\omega_{k_m}=\omega_k,
\]
the eigenvalue alone need not fix a unique direction: different
linearly independent vectors can belong to the same eigenspace.
For a Hermitian operator, an eigenvalue of multiplicity $m$ has
$m$ linearly independent eigenvectors.
\editorial{The source states this conclusion for a repeated
characteristic root without qualifying the operator. In general,
algebraic multiplicity need not equal eigenspace dimension. For
example, $\left[\begin{smallmatrix}\omega&1\\0&\omega\end{smallmatrix}\right]$
has a double root but only one independent eigenvector. The equality
used in the following notes holds for Hermitian operators.}
